\begin{table}[htbp]
\centering
\caption{Empirical Application Results}
\label{tab:application}
\resizebox{\ifdim\width>\linewidth\linewidth\else\width\fi}{!}{% shrinks the table only if it is wider than the text
\begin{tabular}{lccccccccccccccccc}
 &  &  &  &  & \multicolumn{2}{c}{$J_0=20$} & \multicolumn{7}{c}{EKW $p$-values} & \multicolumn{4}{c}{Kudrin (2025) $p$-values} \\
\cmidrule(lr){6-7} \cmidrule(lr){8-14} \cmidrule(lr){15-18}
 & $n$ & Articles & Med.\ $\nu$ & $\Delta_\nu$ & $p$ & $\widehat{B}$ & LCM & Fisher & Disc. & Binomial & CS1 & CS2B & Min & $T_\infty$ & $T_H$ & $T_H^{\mathrm{hybrid}}$ & Min \\
\midrule
RCT & 6503 & 141 & 185 & 0.0014 & .00 & .0032 & .95 & 1.00 & .85 & .16 & .48 & .69 & .16 & .01 & .00 & .00$^{\dagger}$ & .00 \\
IV & 4901 & 277 & 262 & 0.0013 & .00 & .0034 & 1.00 & 1.00 & .14 & .11 & .66 & .24 & .11 & .38 & .05 & .05$^{\dagger}$ & .05 \\
DID & 5142 & 236 & 293.5 & 0.0014 & .01 & .0017 & .99 & 1.00 & .03 & .21 & .13 & .02 & .02 & .07 & .02 & .02$^{\dagger}$ & .02 \\
RDD & 2722 & 82 & $\cdot$ & $\cdot$ & .03 & .0007 & 1.00 & 1.00 & .98 & .66 & .99 & .76 & .66 & .26 & .19 & .19$^{\dagger}$ & .19 \\
\bottomrule
\end{tabular}}

\medskip
\begin{minipage}{\linewidth}
Results for the Brodeur et al. (2020) data. Following Kudrin (2025), we drop tests whose reported coefficient is missing or zero.
As in Appendix D.3, we divide each coefficient by the smallest standard error consistent with the reported one.
We drop $t$-statistics whose standard error is reported as zero or whose rounding unit in $t$-units exceeds 0.3.
Columns headed $J_0$ report $p$-values and breakdown statistics $\widehat{B}$ of the projection test with $J=40$.
Med.\ $\nu$ is the median effective sample size as measured by the median number of clusters in a random sample of 100 articles for each causal inference method.
In ambiguous cases, the smallest cluster count was always used.
$\Delta_\nu=\frac{1}{100}\sum_{k=1}^{100}\sup_h\|\varphi(t-h)-t_{\nu_k,h}(t)\|_T$ is the average difference between the PDFs of the standard normal and Student-$t$ over the sample of 100 effective sample sizes collected (computed numerically).
We form $\mathcal{X}$ using a grid of 3000 elements evenly spaced between $\pm 6.5$ plus the scalar 9999.
Discretization incurs error $\mathcal{E}(L,\delta)=0.00063$, which is already accounted for in $\widehat{B}$ and the $p$-value.
The discretization was chosen to be made of round numbers and result in error relatively small compared to $\Delta_\nu$.
EKW $p$-values are for the tests of Elliott et al. (2022): the least concave majorant test (LCM) and Fisher's test (Fisher) on $p$-values in $(0,0.15]$, the density discontinuity test at 0.05 (Disc.), the binomial test on $[0.04,0.05]$ (Binomial), and the Cox and Shi (2023) histogram-based tests on $(0,0.15]$ with 15 bins for monotonicity (CS1) and for 2-monotonicity with bounds (CS2B).
Kudrin (2025) $p$-values are for his KS-type test ($T_\infty$), histogram-based test ($T_H$) and hybrid histogram-based test ($T_H^{\mathrm{hybrid}}$), using the sinc kernel with $\kappa=3$ and, for $T_H$ and $T_H^{\mathrm{hybrid}}$, the points $0.01,0.02,\ldots,0.15$.
All three cluster by article: $T_H$ and $T_H^{\mathrm{hybrid}}$ use cluster-robust covariance matrices and $T_\infty$ uses an article-level bootstrap.
$^\dagger$The deconvolution estimate already satisfies the hybrid test's shape restrictions, so $T_H^{\mathrm{hybrid}}=T_H$.
Min is the smallest $p$-value in each group of tests.
The projection test and $T_\infty$ use 1000 bootstrap repetitions.
Computation used osqp package in R by Stellato et al. (2024). Data: Brodeur et al. (2020).
Transparency R package available at: \url{https://github.com/sfaridan/When_is_p_hacking_detectable}
\end{minipage}
\end{table}
